\documentclass[11pt,letterpaper]{article}
\usepackage[T1]{fontenc}
\usepackage[utf8]{inputenc}
\usepackage{lmodern}
\usepackage[margin=1in,headheight=14pt]{geometry}
\usepackage{amsmath,amssymb,amsthm,mathtools}
\usepackage{microtype,booktabs,longtable,array,enumitem}
\usepackage{xcolor,listings,xurl,hyperref,fancyhdr}
\hypersetup{colorlinks=true,linkcolor=blue!45!black,citecolor=blue!45!black,
 urlcolor=blue!45!black,pdftitle={Siegel's theorem and intuitionistic PRA},
 pdfauthor={Research report prepared for Vladimir Reshetnikov}}
\pagestyle{fancy}\fancyhf{}
\fancyhead[L]{\small Siegel's theorem and intuitionistic PRA}
\fancyhead[R]{\small Constructive scope and explicit cases}
\fancyfoot[C]{\thepage}
\setlength{\parskip}{0.35em}
\setlength{\parindent}{0pt}
\setlist{itemsep=0.25em,topsep=0.4em}
\setlength{\emergencystretch}{2em}
\newtheorem{theorem}{Theorem}[section]
\newtheorem{lemma}[theorem]{Lemma}
\newtheorem{proposition}[theorem]{Proposition}
\newtheorem{corollary}[theorem]{Corollary}
\theoremstyle{definition}
\newtheorem{definition}[theorem]{Definition}
\newtheorem{remark}[theorem]{Remark}
\newcommand{\ipra}{\mathrm{iPRA}}
\newcommand{\N}{\mathbb N}\newcommand{\Z}{\mathbb Z}
\newcommand{\Q}{\mathbb Q}
\newcommand{\abs}[1]{\lvert #1\rvert}
\DeclareMathOperator{\ord}{ord}
\DeclareMathOperator{\height}{ht}
\lstset{basicstyle=\ttfamily\small,breaklines=true,columns=fullflexible,
 keepspaces=true,frame=single,framerule=0.3pt,rulecolor=\color{black!25},
 showstringspaces=false}
\title{Siegel's Theorem and Intuitionistic Primitive Recursive Arithmetic\\[0.6em]
\large Explicit Runge bounds, complete positive-genus examples,\\
and the unresolved general proof}
\author{Research report prepared for Vladimir Reshetnikov}
\date{15 November 2026}
\begin{document}
\maketitle
\begin{abstract}
The full requested proof of Siegel's theorem in intuitionistic primitive
recursive arithmetic is \emph{not established in this report}. Positive
height-bounded finiteness would have strong effectivity consequences; a
cardinality formulation is weaker and requires a different proof. We explain
both distinctions, inspect the relevant ProveIt foundation, and identify the
unproved general steps rather than postulating them.

The constructive results proved here are unconditional and uniform. A finite
polynomial identity $D^2f=A^2+R$, with $D>0$, $\deg R<\deg A$, and $R\ne0$,
gives the explicit bound
$|x|\le\sum_{i<\deg A}|a_i|+\sum_i|r_i|$ for every integral solution of
$y^2=f(x)$. Such an identity is constructed by primitive recursive rational
arithmetic whenever $f$ is nonsquare, of positive even degree, and has nonzero
square leading coefficient. Finite enumeration is then exhaustive. In
particular, $y^2=x^{2m}+x+1$ has exactly the four points
$(-1,\pm1),(0,\pm1)$ for every $m\ge2$; these curves have genus $m-1$.
We also prove a sharper quartic bound and a complete six-point example.
The archive contains the detailed proofs, exact reference algorithms,
compiled native Lean arithmetic and search lemmas, and a dependency-ordered
plan for embedded iPRA derivations. The general Siegel and polynomial-list
targets in the Lean specifications remain unproved.
\end{abstract}
\tableofcontents
\newpage
\section{The requested theorem and the actual result}

\subsection{The mathematical statement}

Fix a number field $K$, its ring of integers $\mathcal O_K$, and a smooth
geometrically integral affine curve $C$ over $K$ with a specified embedding
in affine space. An integral point is a tuple in $\mathcal O_K^n$ satisfying
the defining equations of that affine presentation. Let $\overline C$ be
the smooth projective completion and let $g$ be its geometric genus.
Siegel's theorem states that
\[
 g\ge1\quad\Longrightarrow\quad C(\mathcal O_K)\text{ is finite}.
\]
The more general Siegel--Mahler statement permits a finite set $S$ of places
containing the infinite places and the ring $\mathcal O_{K,S}$ of elements
integral outside $S$. If $r$ is the number of geometric points in
$\overline C\setminus C$, the sufficient condition is $g\ge1$ or $r\ge3$.
An affine curve has $r\ge1$, so this condition is equivalent to
$2g-2+r>0$. We retain the ordinary-integral, positive-genus statement as
the requested goal. The classical background is discussed in
\cite{BiluIllengo2009,CorvajaZannier2002}.

The embedding and the deleted points matter. A birational equivalence of
smooth projective curves can introduce denominators, and it generally does
not take the original integral-point problem to the integral-point problem
of an arbitrary new affine equation. Thus replacing a curve by a convenient
equation requires a proved, effective comparison of the integral models.

For example, the existence of hyperelliptic equations for genus-two curves
does not settle all their integral-point problems. Corvaja--Lombardo--Zannier
explain precisely this distinction for curves with one point removed, and
report that a general terminating procedure for their integral-point sets
is not known \cite[Introduction and Remark 1.2]{CorvajaLombardoZannier2024}.
Our explicit results concern their stated equations over $\Z$. They do not
erase this dependence on the affine presentation.

\subsection{Why a full affirmative answer is not claimed}

The usual classical theorem does not, in general, provide a computable bound
on the sizes of the points. A constructive theorem producing a finite
complete list, or a bound followed by finite search, would supply that
missing information. In iPRA the witness consequence is particularly
strong: under decidable input hypotheses a uniform proof would yield a
primitive recursive bound. Section~\ref{sec:logical-audit} proves this
conditional implication carefully.

This is a serious obstacle, not a proof of impossibility. We neither prove
that Siegel's theorem is independent of iPRA nor claim that the effective
problem cannot be solved. Moreover, the user did not specify that
``finite'' must mean a positive complete list. A bound on the \emph{number}
of points is a different meaningful constructive target. That target does
not automatically settle the height problem. A complete iPRA proof of the
general cardinality formulation is also absent here.

The recent primary source above and the older effective results describe
many cases with computable solutions. This report does not claim that
positive genus itself forces ineffectivity, or that all genus-two models
are beyond effective methods. The open general step and the elementary
families below must be distinguished.

\subsection{Results established in the report}

The main arithmetic theorem applies to a checkable finite certificate,
not to a promised numerical bound. Write $A(X)=a_mX^m+\cdots+a_0$ with
$a_m\ne0$, and let $R(X)=\sum_i r_iX^i$. Suppose
\[
 D^2f(X)=A(X)^2+R(X),\qquad
 D\in\Z_{>0},\quad m\ge1,\quad\deg R<m,\quad R\ne0.
\]
Then every integral point $y^2=f(x)$ satisfies
\[
 |x|\le B_x:=\sum_{i<m}|a_i|+\sum_i|r_i|,
 \qquad
 |y|\le B_y:=\sum_i |f_i|\max(1,B_x)^i.
\]
These are explicit integer bounds, and both the certificate and the
enumeration can be computed by primitive recursion. The hypotheses are
verified coefficient by coefficient. In particular, the certificate
does not contain an assertion about all solutions as an unchecked field.

We give a uniform certificate constructor for every nonsquare integral
polynomial $f$ of even positive degree whose leading coefficient is a
nonzero integer square. When $f$ is squarefree of degree $2m\ge4$, this
gives a genuine positive-genus instance of Siegel finiteness. The proof
is an elementary instance of Runge's method: two integral quantities
cannot be nonzero factors of a smaller integer.

We additionally give a sharper factor-divisor algorithm for
$y^2=x^4+ax^2+b$, a complete six-point classification for $a=7,b=1$,
and a four-point classification valid in every positive genus. No
claim of novelty is made for the underlying factorization or Runge
principle; the purpose is an explicit arithmetic proof and a transparent
formalization interface.

\subsection{Repository snapshot and reuse}

The repository was inspected through the GitHub connector at revision
\begin{center}\small\ttfamily
7eeb3c0e821594d5d5ebcfbaf4bdba64819e3502.
\end{center}
The actual foundation supplies primitive recursive syntax, intuitionistic
natural deduction, quantifier-free induction, and faithful translation
to the concise native calculus. We use those interfaces in
Section~\ref{sec:lean}; we do not infer an arithmetic theorem from
bi-interpretability alone.

The repository's proof-packet index explicitly says that the packets were
filed as delivered, without being rerun or verified at placement.
The inspected Faltings package explicitly leaves general Faltings in iPRA
unproved; the Mordell--Weil package retains a covering-list hypothesis in
its general descent theorem; and the Roth package leaves a general
effective cutoff unproved. Their existence therefore supplies no verified
general Siegel height bound. These observations are limited to the
identified files and are not an exhaustive correctness review of ProveIt.
Exact paths and the pinned source manifest accompany the archive
\cite{ProveIt2026}.

\subsection{Status at a glance}

\begin{center}
\begin{tabular}{@{}p{0.48\linewidth}p{0.43\linewidth}@{}}
\toprule
Claim & Status in this archive\\
\midrule
General positive Siegel finiteness in iPRA & Not proved\\
General cardinality Siegel finiteness in iPRA & Not proved\\
General negative translation in iPRA & Not proved\\
Explicit Runge bound and certificate constructor & Fully developed arithmetic proof\\
Exhaustive search for the certified family & Proved; executable reference code\\
Every-genus four-point family & Complete arithmetic classification\\
Biquadratic factor bound and six-point example & Written proof and compiled native Lean proof\\
Embedded \texttt{Prf PRAAx} derivations & Detailed implementation plan; not constructed\\
\bottomrule
\end{tabular}
\end{center}

% Standalone section for inclusion in the main report.
% Requires amsmath, amssymb, and bibliography keys Salehi2002, Wehmeier1997.
\section{What constructive finiteness would assert}
\label{sec:logical-audit}

The word ``finite'' requires particular care in this problem.  Several
classically equivalent statements about integral points have different
intuitionistic meanings.  A proof of one of them must not silently be advertised
as a proof of another.  In particular, a bound on the \emph{number} of integral
points need not locate any point, whereas a bound on their \emph{height}
immediately gives a complete finite search.  The distinction is relevant before
one asks which induction principles a proof uses.

\subsection{A finite arithmetic interface}

For this discussion, work over $\mathbb Q$ and use a fixed affine presentation.
An input $c$ encodes an integer $m\geq 1$ and finitely many polynomials
$F_1,\ldots,F_r\in\mathbb Z[X_1,\ldots,X_m]$.  It may also contain a finite
certificate $w$ of additional hypotheses.  Signed integers, finite sequences,
and polynomial coefficient arrays are represented by natural-number codes.
Use the concrete signed-integer encoding
\[
 \operatorname{enc}(z)=
 \begin{cases}2z&z\geq0,\\-2z-1&z<0.\end{cases}
\]
Its inverse, absolute value, the elementary ring operations, polynomial
evaluation, sequence length, and sequence entries are primitive recursive.
Choose an injective canonical coding of finite tuples, so each tuple has
exactly one accepted code.  The point predicate rejects noncanonical codes;
equality and list membership therefore compare points rather than multiple
representations of one point.  Malformed codes receive a harmless default
interpretation in the arithmetic operations, so these operations are total
functions on natural numbers.

Let $P(c,p)$ mean that $p$ is a well-formed $m$-tuple of integers and
$F_i(p)=0$ for every $i\leq r$.  Let
\[
 h(c,p)=\max_{1\leq j\leq m}|p_j|
\]
on well-formed tuples.  Both $P$ and $h$ are primitive recursive.  Here $h$ is
an affine coordinate size; no theory of canonical heights is needed.  Finally,
let $G(c)$ be a primitive recursive test for the supplied admissibility
certificate.  We do \emph{not} assume that every desired geometric property has
already been compiled into such a test.  That compilation, its soundness, and
its coverage are separate obligations.

This interface isolates the logical issue without assuming a constructive
theory of algebraic curves.  It also permits statements conditional on a
specific, independently verified certificate.  If a later development chooses
a decidable geometric presentation directly, $G$ can instead be its validity
test.  Conversely, an unbounded, nondecidable hypothesis must not be replaced
by $G$ without justification.

Write
\[
 \operatorname{Bound}(c,B):\equiv
 \forall p\bigl(P(c,p)\longrightarrow h(c,p)\leq B\bigr).
 \tag{\textsf{H}}
\]
The quantifier over $p$ is unbounded.  Consequently,
$\operatorname{Bound}(c,B)$ is generally a $\Pi^0_1$ formula, although every
individual polynomial evaluation and comparison is decidable.  Decidability
of those individual tests does not make the universal assertion decidable.

\subsection{Height bounds and complete lists}

A positive finite-list statement has the form
\[
 \exists L\ \forall p\bigl(P(c,p)\longleftrightarrow p\in L\bigr),
 \tag{\textsf{L}}
\]
where $L$ is a natural-number code for a finite sequence.  Repetitions may be
removed by a primitive recursive filter.  Thus the sequence formulation is
equivalent to a coded duplicate-free list, without invoking arbitrary sets.

There are direct constructive transformations between $(\mathsf L)$ and
$\exists B\,\operatorname{Bound}(c,B)$.  Given $L$, compute the maximum of the
heights of its entries, with maximum $0$ for the empty list.  Every point
occurs in the list, hence its height is at most this maximum.  This direction
only needs the covering implication $P(c,p)\to p\in L$; if the list contains
extraneous tuples, filtering by $P$ removes them.

For the converse, suppose a particular $B$ and a proof of
$\operatorname{Bound}(c,B)$ are available.  Enumerate the box
\[
 \{-B,-B+1,\ldots,B\}^{m}.
\]
It has exactly $(2B+1)^m$ entries.  For example, digits in base $2B+1$
enumerate coordinate tuples, with each digit $d$ interpreted as $d-B$.
Evaluate the defining polynomials on each tuple and retain the tuples for
which every value is zero.  All loops have explicit finite limits.  Their
composition is primitive recursive, including when $m$ and $B$ vary.
Every retained tuple is a point by the filter test.  Every point lies in the
box by the given bound and is therefore retained.  This proves completeness.

The associated iPRA argument uses only elementary coding and induction on
finite loops.  A statement that a loop has inspected its first $k$ entries is
bounded; in the language with all primitive recursive functions it can be
expressed by a quantifier-free characteristic equation.  Thus the usual
induction proof of filter correctness does not require induction over an
arbitrary unbounded formula.  This observation justifies the finite search
\emph{after} a height bound has been proved; it does not supply that bound.

For integral points over a number field, one must additionally specify a
finite encoding of the field and integral elements and give an effective
finite enumeration below the chosen height.  Those additions do not remove
the distinction between bounding the number of points and bounding their
locations.  The rational case already exhibits the logical issue.

\subsection{A bound on the number of points is weaker}

Define $\operatorname{AtMost}(c,N)$ by saying that every list of $N+1$ points
contains a repeated entry.  More explicitly,
\[
 \forall s\Bigl(
   \operatorname{length}(s)=N+1\ \wedge\
   \operatorname{AllPoints}(c,s)
   \ \longrightarrow\
   \operatorname{HasRepeat}(s)\Bigr).
 \tag{\textsf{C}}
\]
The two sequence predicates are primitive recursive: their internal
quantification is bounded by the sequence length.  This is therefore also a
$\Pi^0_1$ assertion.  A complete list gives $(\mathsf C)$ with its length as
$N$, by the finite pigeonhole principle.  A height bound gives it with
$N=(2B+1)^m$.

The reverse implication is not intuitionistically valid in general, even
for primitive recursive predicates.  To see the reason, let $R(e,t)$ mean
that a deterministic program with code $e$, on a fixed input, halts
\emph{for the first time} after exactly $t$ steps.  Bounded simulation makes
$R$ primitive recursive.  Finite computation and determinism prove that
there cannot be two different such times.  In particular the numerical
bound $1$ works uniformly for the number of witnesses of $R(e,\cdot)$.

Nevertheless, no total computable function $b(e)$ can bound every such
halting time.  If it did, simulate program $e$ for $b(e)$ steps.  If it has
halted, answer yes; otherwise the bound guarantees that it never halts, so
answer no.  This would decide the halting problem.  Thus a uniform count
bound, even the constant $1$, does not in general give a computable height
bound or a complete list.

This example concerns arbitrary decidable subsets of the natural numbers.
It is \emph{not} a reduction of the halting problem to integral points on
the curves under discussion.  It supplies no undecidability result for
Siegel's theorem.  Its purpose is narrower: it rules out a purely logical
passage from a cardinality bound to an exhaustive search bound.

Knowing an upper bound $N$ also does not justify searching until $N$ points
have been found.  There may be fewer than $N$ points, or none.  Such a
search would then fail to stop.  A proof providing $N$ actual distinct
points together with an upper bound $N$ is different: the provided points
form a complete list.  The positive witnesses in the lower-bound statement
are essential.

Accordingly, the uniform assertion
\[
 \forall c\bigl(G(c)\to\exists N\,
                         \operatorname{AtMost}(c,N)\bigr)
 \tag{\textsf{CountSiegel}}
\]
would extract a primitive recursive bound on the \emph{number} of points.
It need not solve the effective problem of locating all points.  It would
still require a complete proof in the specified arithmetic.  The existence
of classical quantitative results is not, by itself, such a proof.  The
phrase ``finitely many'' in the requested theorem therefore leaves a real
choice of constructive formulation; this report does not identify
$(\mathsf{CountSiegel})$ with the stronger assertion below.

\subsection{The precise witness-extraction consequence}

For the usual presentation of iPRA, the relevant metatheorem is stronger
than mere recursive realizability.  If
\[
 \mathrm{iPRA}\vdash\forall u\exists v\,A(u,v),
\]
where $A$ is an arbitrary arithmetical formula, then there is a primitive
recursive function symbol $f$ such that
\[
 \mathrm{iPRA}\vdash\forall u\,A(u,f(u)).
 \tag{\textsf{W}}
\]
Salehi records this as Fact~3.5, with a reference to Wehmeier's
primitive-recursive choice theorem
\cite{Salehi2002,Wehmeier1997}.  Multiple input variables are combined by
primitive recursive pairing.  In particular, the theorem is not restricted
to quantifier-free $A$; its applicability to a matrix containing an
unbounded universal quantifier is exactly what is needed here.

Here is the consequence for a hypothetical proof of positive Siegel
finiteness.  Suppose
\[
 \mathrm{iPRA}\vdash
 \forall c\bigl(G(c)\longrightarrow
          \exists B\,\operatorname{Bound}(c,B)\bigr),
 \tag{\textsf{PositiveSiegel}}
\]
and suppose iPRA proves the decidability of $G$.  For any fixed $c$, reason
by that decidability.  If $G(c)$, use $(\mathsf{PositiveSiegel})$ and retain
its witness $B$.  If $\neg G(c)$, take $B=0$; the implication from $G(c)$
is vacuous.  These two cases prove, intuitionistically,
\[
 \mathrm{iPRA}\vdash\forall c\exists B
       \bigl(G(c)\longrightarrow\operatorname{Bound}(c,B)\bigr).
\]
Apply $(\mathsf W)$ with
$A(c,B)\equiv G(c)\to\operatorname{Bound}(c,B)$.  It supplies a primitive
recursive $b$ and an iPRA proof of
\[
 \forall c\bigl(G(c)\longrightarrow
                         \operatorname{Bound}(c,b(c))\bigr).
 \tag{\textsf{ExtractedBound}}
\]
No unrestricted independence-of-premise principle has been used.  The
case distinction on the explicitly decidable certificate predicate is
what permitted moving the existential quantifier.

Composing $b$ with the box enumeration and polynomial filter gives a
primitive recursive procedure which outputs a complete list on every
certified admissible input.  It also decides whether that input has any
integral points, by checking whether the resulting list is empty.  The
numerical bound need not be remotely practical; primitive recursiveness
is a structural conclusion about its construction, not a claim of
feasible running time.

The same argument applies to $(\mathsf{CountSiegel})$, but with
$A(c,N)=G(c)\to\operatorname{AtMost}(c,N)$.  Its extracted function only
supplies $N$.  Substituting this function into a height-bounded search
would be an error: no height bound occurs in the conclusion of the
witness theorem.

There are two further qualifications.  First, when $c$ contains a
certificate, the extracted bound is allowed to depend on that certificate.
Obtaining an algorithm from raw equations requires a procedure that
produces a valid certificate for every admissible raw input.  A
primitive recursive certificate checker alone does not supply such a
procedure.  Second, if admissibility is instead represented by a
nondecidable formula, a realizability argument may require a realizer of
that hypothesis.  It is not legitimate to discard that additional input
and announce a bound depending only on the curve's coefficients.

\subsection{Fixed instances and external schemas}

A proof for one specified curve and a uniform theorem over curve codes
have different strengths.  If iPRA proves
$\exists B\,\operatorname{Bound}(\bar c,B)$ for one fixed numeral
$\bar c$, numerical existence provides a particular standard integer
$b$ with an iPRA proof of $\operatorname{Bound}(\bar c,\bar b)$.
There is no conclusion that the same bound works for other curves.
For that particular curve, however, the proof supplies precisely the
information required for a complete finite search.

There is a subtle intermediate case.  Suppose one asserts externally
that \emph{every} admissible standard curve code $c$ has some iPRA proof
of $\exists B\,\operatorname{Bound}(\bar c,B)$, without presenting a
uniform proof of $(\mathsf{PositiveSiegel})$.  This assertion still has
an effectivity consequence.  On a given admissible input, enumerate iPRA
proofs.  Numerical existence guarantees that eventually one finds a
proof of $\operatorname{Bound}(\bar c,\bar b)$ for some numeral
$\bar b$.  Return $b$ and perform the finite search.  Proof enumeration
and recognition are effective, so this is a partial recursive algorithm
which terminates on every admissible input covered by the asserted
schema.  Its running time, and the resulting bound as a function of $c$,
need not be primitive recursive.

Thus a nonuniform assertion of \emph{provability of all instances}
must not be confused with bare classical existence of a bound for
each instance.  The latter does not furnish recognizable iPRA proofs
to enumerate.  Nor does a collection of proofs for selected examples
establish that all admissible instances are covered.  These distinctions
prevent both an exaggerated impossibility claim and an understated
claim about what a proposed constructive theorem would accomplish.

\subsection{Negations, translations, and hidden induction}

A possible weaker target is
\[
 \forall c\bigl(G(c)\longrightarrow
              \neg\neg\exists B\,\operatorname{Bound}(c,B)\bigr).
 \tag{\textsf{NegativeSiegel}}
\]
The witness theorem does not remove the double negation.  In particular,
ordinary Markov-style reasoning for a decidable search predicate does
not apply directly: the predicate being searched for is
$\operatorname{Bound}(c,B)$, an unbounded universal assertion.  Checking
a large finite region cannot verify it.

The statement that points are not unbounded,
\[
 \neg\forall B\exists p\,
                    \bigl(P(c,p)\wedge h(c,p)>B\bigr),
\]
is another possible negative formulation.  It follows from positive
boundedness and also from its double negation.  One should not assume
that all these negative formulations are intuitionistically equivalent;
interchanging negation and unbounded quantifiers itself requires
justification.  Likewise, the absence of a specified infinite sequence
of distinct points is not a coded complete list of the points.

Negative translation is useful only with the correct source theory.
A classical \emph{PRA} proof can be translated into an iPRA proof of
its negative translation, because the relevant defining equations and
quantifier-free induction are preserved.  A proof in arbitrary classical
mathematics does not automatically become an iPRA proof by the same
operation.  Its induction, choice, comprehension, and algebraic
existence principles must first be accounted for.  Therefore a classical
proof of Siegel's theorem cannot simply be relabeled a proof of
$(\mathsf{NegativeSiegel})$ in iPRA.

Similar care is needed with formalization.  A Lean proof about
$\mathbb Z$, polynomials, and finite lists can establish the correctness
of a finite search.  A Lean theorem whose hypothesis already asserts a
height bound is a conditional result.  Neither one, by itself, proves
that iPRA derives the required height bound.  For the latter claim,
one needs an explicit translation to object-language arithmetic and a
derivation whose induction formulas satisfy the stated restriction.
Using Lean's unrestricted natural-number induction in the metatheory
does not discharge that obligation.

\subsection{A finite kernel that can be completed independently}

Several useful lemmas do not depend on the unresolved global argument.
First, primitive recursive evaluation of a polynomial code agrees with
the corresponding finite sum of monomials.  A proof proceeds through
the coefficient list and then through each exponent vector; all
intermediate identities are ring equations.  Second, decoding the
base-$(2B+1)$ expansion of an index below $(2B+1)^m$ gives exactly one
tuple in the box.  Conversely, encoding the shifted coordinates of
any box tuple returns such an index.  Both assertions reduce to
bounded division and remainder identities.

Third, filtering a finite sequence by the point predicate is sound,
and every original entry satisfying that predicate occurs in the
filtered sequence.  Fourth, the maximum-height operation bounds the
height of each list entry.  Finally, the pigeonhole argument converts
a complete list into a bound on the number of distinct points.  These
proofs may be organized with characteristic functions for bounded
sequence predicates, so their loop induction remains quantifier-free.
Their correctness is uniform in the coefficient data, the number of
coordinates, and the supplied numerical bound.

These results constitute a genuine computational kernel, but its
interface must retain the hypothesis
$\operatorname{Bound}(c,B)$.  An executable routine that receives $B$
and returns the filtered box list is already useful; calling its
output ``all integral points'' without verifying that hypothesis is
not justified.  Likewise, checking finitely many examples of the
routine checks implementation behavior, not the global theorem.

An arithmetic proof checker can instead check a finite derivation of
the universal bound.  This is a stronger certificate than merely
testing points up to a chosen cutoff.  Its syntax, inference rules,
and induction restriction must be specified and its soundness
established in the metatheory.  Having such a checker does not create
the required derivation.  Separating the finite search kernel, the
proof checker, and the mathematical height-bound argument prevents a
correct piece of infrastructure from concealing the missing theorem.

\subsection{What this audit establishes}

The established logical conclusion is conditional and precise: a
uniform iPRA proof of positive, height-bounded Siegel finiteness on a
decidably certified family yields a primitive recursive height bound
and complete enumeration on that family.  A uniform iPRA proof of
cardinality finiteness yields only a primitive recursive cardinality
bound.  An admissible negative formulation has still different
computational content.

None of these observations proves that Siegel's theorem, in any of
these formulations, is unprovable in iPRA.  They identify what must be
proved and what cannot be supplied by a purely formal change of
language.  In particular, stating that the positive version would
settle an unresolved effective problem is not an independence theorem.
A completed contribution would have to provide the missing arithmetic
argument, or explicitly choose and prove a weaker constructive
formulation.  The finite coding and extraction results above remain
useful in either direction, while making the remaining mathematical
work visible.

\section{A complete constructive subcase: a finite Runge certificate}
\label{sec:runge-certificate}

The following result is an actual, uniform arithmetic theorem.  Its input is
a finite algebraic certificate, its output is an explicit box containing every
integral point, and its proof uses only integer arithmetic and bounded sums.
It supplies a substantial effective subcase of the requested theorem, including
smooth curves of arbitrarily large genus.  It does not supply the missing
certificate for an arbitrary curve to which Siegel's theorem applies.

Throughout this section a polynomial is a finite coefficient list, in ascending
order.  Trailing zero coefficients are removed; the empty list represents zero.
Its degree, coefficients, sum of absolute coefficients, products, and value at
an integer are primitive recursive operations.  Polynomial equalities below
are therefore finite, decidable coefficient identities.  Integers can be
represented by any fixed primitive recursive signed coding.  Statements using
absolute values are abbreviations for statements about the natural-number
absolute-value function on those codes.

\subsection{Two elementary bounds}

\begin{lemma}[An integer root bound]
\label{lem:runge-root-bound}
Let $G(X)=\sum_{i=0}^{r}g_iX^i$ be a nonzero integer polynomial.  If $r=0$,
it has no root.  If $r\geq1$ and $G(x)=0$ for an integer $x$, then
\[
 |x|\leq \sum_{i<r}|g_i|.
\]
\end{lemma}
\begin{proof}
The assertion is immediate for $x=0$.  Otherwise put $t=|x|\geq1$ and
$L=\sum_{i<r}|g_i|$.  Since the nonzero integer $g_r$ has $|g_r|\geq1$,
the equation $g_rx^r=-\sum_{i<r}g_ix^i$ gives
\[
 t^r\leq |g_r|t^r
   \leq\sum_{i<r}|g_i|t^i
   \leq L t^{r-1}.
\]
Cancellation of the positive integer $t^{r-1}$ gives $t\leq L$.  The constant
case says simply that a nonzero integer is not zero.  All sums are finite;
the comparison $t^i\leq t^{r-1}$ is proved by induction on the exponent.
\end{proof}

\begin{lemma}[A nonzero product]
\label{lem:runge-product-bound}
If $u,v$ are integers with $uv\neq0$, then
\[
 |uv|\geq\max\{|u|,|v|\},\qquad
 |u|+|v|\leq |uv|+1.
\]
Moreover, if $v-u=2z$, then $\max\{|u|,|v|\}\geq |z|$.
\end{lemma}
\begin{proof}
Write $U=|u|$ and $V=|v|$.  Both are positive integers, so $UV\geq U,V$.
Expanding $(U-1)(V-1)\geq0$ gives $U+V\leq UV+1$.
Finally,
$2|z|=|v-u|\leq U+V\leq2\max(U,V)$, and cancellation of $2$ proves the
last assertion.  These are finite inequalities, with no choice of a prime
factor and no factorization theorem needed.
\end{proof}

\subsection{The certificate theorem and its complete proof}

\begin{theorem}[Elementary effective Runge certificate]
\label{thm:runge-certificate}
Let $f,A,R\in\mathbb Z[X]$ and $D\geq1$ satisfy the coefficient identity
\[
                 D^2 f=A^2+R.
\tag{RC}
\]
Assume that $m=\deg A\geq1$, that $R\neq0$, and that $\deg R<m$.
Write
\[
 \alpha=\sum_{i<m}|a_i|,\qquad
 \beta=\sum_i|r_i|,\qquad H=\alpha+\beta,
 \qquad K=\sum_i|f_i|H^i.
\]
Then every integer solution of $y^2=f(x)$ satisfies
\[
                         |x|\leq H,
                    \qquad |y|\leq K.
\]
Consequently the integral points have a primitive recursive complete
enumeration, uniformly in $f,D,A,R$ satisfying the displayed conditions.
\end{theorem}

\begin{proof}
Because $R$ is nonzero, $\beta\geq1$, hence $H\geq1$.
Let $y^2=f(x)$ and put $t=|x|$.  Integer order is decidable.  It suffices to
show that its alternative $t>H$ leads to a contradiction.

First, $R(x)\neq0$.  This is immediate if $R$ is constant.  If its degree is
$r\geq1$, Lemma~\ref{lem:runge-root-bound} would otherwise give
$t\leq\sum_{i<r}|r_i|\leq\beta\leq H$.

Second, $A(x)$ is larger in absolute value than $R(x)$.  Since $t\geq1$ and
the leading coefficient of $A$ is a nonzero integer, the reverse triangle
inequality gives
\begin{align*}
 |A(x)|
 &\geq |a_m|t^m-\sum_{i<m}|a_i|t^i\\
 &\geq t^m-\alpha t^{m-1}
   =(t-\alpha)t^{m-1}
   >\beta t^{m-1}.
\end{align*}
The strict inequality follows from $t>\alpha+\beta$ and
$t^{m-1}\geq1$.  Every exponent occurring in $R$ is at most $m-1$, so
\[
                  |R(x)|\leq\beta t^{m-1}<|A(x)|.
\]

Now use (RC) at the particular integer $x$ and substitute $y^2=f(x)$:
\[
              (Dy-A(x))(Dy+A(x))=R(x).
\]
The two factors are nonzero integers.  Their difference is $2A(x)$, so
Lemma~\ref{lem:runge-product-bound} gives
\[
 |R(x)|\geq
 \max\{|Dy-A(x)|,|Dy+A(x)|\}\geq |A(x)|.
\]
This contradicts the preceding strict inequality.  We have proved $t\leq H$.

For the second coordinate,
\[
 y^2=f(x)\leq |f(x)|\leq\sum_i|f_i|\,|x|^i
                         \leq\sum_i|f_i|H^i=K.
\]
For every integer $y$, the natural number $n=|y|$ satisfies $n\leq n^2$:
separate $n=0$ from $n\geq1$ and multiply $1\leq n$ by $n$.
Thus $|y|\leq y^2\leq K$.

Finally, run two bounded loops through $-H\leq x\leq H$ and
$-K\leq y\leq K$, retaining exactly the pairs for which the decidable
identity $y^2=f(x)$ holds.  Every retained pair is a solution by the test.
Every solution is retained by the two bounds just proved.  This produces a
finite list with both soundness and completeness, rather than merely a
numerical bound on the number of possible points.
\end{proof}

The use of a contradiction in this proof is intuitionistically harmless.
We did not infer an existential witness from the negation of its negation.
We computed $H$ and $K$ before considering a point, and used the decidable
integer dichotomy $t\leq H$ or $t>H$ to verify their universal correctness.
Similarly, the conclusion $R(x)\neq0$ follows by excluding equality in a
decidable integer calculation.  No principle for choosing from an infinite
family of finite sets occurs.

The factorization step explains the scope of the method.  The certificate
forces $Dy$ to be close to either $A(x)$ or $-A(x)$ at large $|x|$.  The two
integer factors cannot both remain nonzero while their product has lower
growth than their difference.  An arbitrary affine curve need not admit
this particular description; the theorem contains no algorithm for
manufacturing it in that generality.

The domain of the coordinates is equally essential.  The proof uses the
fact that the absolute value of a nonzero \emph{integer} is at least one.
For rational factors this fails: $u$ may be an arbitrarily small positive
fraction while $v$ is large.  Thus the same displayed polynomial identity
does not prove finiteness of rational points.  Nor does it automatically
prove an analogous assertion for coordinates whose denominators are
supported on a prescribed set of primes.  In that setting the product
argument would have to control several absolute values simultaneously;
the ordinary real absolute value alone no longer gives the required
lower bound for each factor.  The theorem here is deliberately stated
over $\mathbb Z$ with its fixed affine coordinates.

This distinction also explains why a birational change of variables must
not be inserted silently.  A rational transformation can introduce
denominators and can omit exceptional points.  To transfer the explicit
box to a different equation, one would have to give the transformation,
prove which integral points it covers, bound any denominators it creates,
and handle all exceptional fibers.  None of those obligations is needed
for the displayed families, whose integral points are bounded directly in
their original coordinates.  A polynomial identity checked on the actual
input equation is therefore part of the certificate, rather than an
informal assertion that the curve resembles an already-solved example.

\subsection{Constructing the certificate by completing a polynomial square}

The certificate is not an oracle for the following broad family.  Suppose
$f$ has degree $2m\geq2$ and its leading coefficient is a positive integer
square $s^2$, where $s\geq1$.  Construct a rational polynomial
\[
                      Q(X)=sX^m+\sum_{j<m}q_jX^j
\]
so that all coefficients of $f-Q^2$ in degrees $m,m+1,\ldots,2m$ vanish.
Start with $q_m=s$.  For $j=m-1,m-2,\ldots,0$, set
\[
 q_j=\frac{1}{2s}
 \left(f_{m+j}-
   \sum_{\substack{j<i<m,\ j<k<m\\i+k=m+j}}q_iq_k\right).
\tag{SQ}
\]
The sum uses only coefficients already defined.  Indeed, a term of degree
$m+j$ in $Q^2$ either uses one coefficient of degree $m$ and the coefficient
of degree $j$, or has both indices strictly greater than $j$.  The former
two terms total $2s q_j$.  Thus (SQ) sets the coefficient in degree $m+j$
equal to $f_{m+j}$ and leaves every previously matched higher coefficient
unchanged.  A finite downward induction verifies
\[
                         \deg(f-Q^2)<m
\]
unless $f-Q^2=0$, when the same coefficient assertion holds.

Take any positive common multiple $D$ of the denominators of the $q_j$.
Their product is sufficient; taking their least common multiple improves
the numerical size.  Define
\[
                        A=DQ,\qquad R=D^2f-A^2.
\]
These are integer polynomials and satisfy (RC).  If $R\neq0$, its degree is
less than $m$, so Theorem~\ref{thm:runge-certificate} applies immediately.
If $R=0$, the constructor reports the exact exceptional case
$f=Q^2$ in $\mathbb Q[X]$ and makes no finiteness claim.  Conversely, if
$f$ is a rational polynomial square, its square root with leading
coefficient $s$ has to obey the same triangular recursion, so this test
does detect that case.  In particular every nonsquare monic integer
polynomial of positive even degree is covered.  Every squarefree such
polynomial is nonsquare and therefore covered as well.

This construction uses exactly $m$ rational coefficient steps.  Each step
uses a bounded sum and rational addition, multiplication, and division by
the explicitly nonzero integer $2s$.  Denominators can be represented
without reduction, or reduced by the primitive recursive Euclidean
algorithm.  Clearing denominators, checking that the remainder is nonzero,
and checking the final coefficient identity are finite arithmetic
operations.  There is no search for a height bound with an unknown stopping
criterion.  The bounds can be large, and no claim of practical efficiency
is part of the theorem.

The restrictions should be kept visible.  For example, the displayed
constructor does not handle an odd-degree cubic or an even-degree
polynomial with a nonsquare positive leading coefficient.  Such a
rejection says that this constructor does not provide the certificate;
it says nothing about whether another method proves finiteness for the
curve.  Conversely, a zero remainder must not be ignored: the integer
polynomial $f(X)=(X^2+1)^2$ has the infinitely many points
$(x,\pm(x^2+1))$.

\subsection{A complete family in every positive hyperelliptic genus}

\begin{theorem}[Uniform exact classification]
\label{thm:runge-all-genera}
For every integer $m\geq2$, the integral solutions of
\[
                         y^2=x^{2m}+x+1
\]
are exactly
\[
                    (-1,-1),\quad(-1,1),\quad(0,-1),\quad(0,1).
\]
\end{theorem}

\begin{proof}
Use $D=1$, $A=X^m$, and $R=X+1$.  Then
$D^2f=A^2+R$, $\deg R=1<m$, $\alpha=0$, and $\beta=2$.
The certificate theorem gives $|x|\leq2$.  It remains to check just five
integers.

For $x=-1$ or $x=0$, the right side is $1$, giving exactly $y=\pm1$.
For $x=1$ it is $3$, strictly between the consecutive squares $1$ and $4$.
For $x=2$, put $k=2^m\geq4$.  The right side is $k^2+3$, and
\[
                         k^2<k^2+3<(k+1)^2,
\]
because $(k+1)^2-k^2=2k+1\geq9$.  For $x=-2$ it is $k^2-1$, and
\[
                        (k-1)^2<k^2-1<k^2,
\]
because $k^2-1-(k-1)^2=2k-2>0$.
An integer square cannot lie strictly between consecutive natural-number
squares: its absolute square root would have to be a natural number
strictly between the two consecutive integers.  This eliminates the
remaining values of $x$.  Substitution checks each of the four listed
points, completing both directions.
\end{proof}

Here the constant list of four points is already a uniform constructive
output; no growth of the list with the genus parameter is necessary.
The exponent $2m$ is part of the input, and exponentiation is primitive
recursive.  The argument uses only the finite coefficient certificate and
inequalities between integers.  It does not assume a general theorem
about rational points or integral points on high-genus curves.

For completeness, these examples really do represent smooth positive-genus
curves, rather than a collection of singular rational equations.  Set
$f_m=X^{2m}+X+1$.  A common zero $\xi$ of $f_m$ and its derivative would
satisfy the polynomial identity
\[
       2m f_m(X)-X f_m'(X)=(2m-1)X+2m
\]
and hence $\xi=-2m/(2m-1)$.  This is a reduced nonintegral rational number
for $m\geq2$.  A rational root of a monic integer polynomial is an integer:
writing it as $p/q$ in lowest terms and clearing denominators shows
$q\mid p^{\deg f}$, so $q=1$.  Thus no such common zero exists, and $f_m$
is squarefree.  An alternative finite certificate is reduction modulo
$2$, where $f_m'$ becomes the constant polynomial $1$.

The affine equation is nonsingular in characteristic zero: a singular
point would have $2y=0$, $f_m(x)=0$, and $f_m'(x)=0$, which has just been
excluded.  At infinity put $t=1/x$ and $w=y/x^m$.  The equation becomes
\[
                         w^2=1+t^{2m-1}+t^{2m}.
\]
At $t=0$ there are exactly the two rational points $w=1$ and $w=-1$;
both are smooth because the derivative with respect to $w$ is $2w\neq0$.
The smooth completion is the double cover of the projective line branched
at the $2m$ simple finite roots of $f_m$, with no ramification at infinity.
The usual hyperelliptic genus calculation is
$2g-2=2(-2)+2m$, giving $g=m-1$.  This geometric identification explains
the range of the examples.  The arithmetic classification above is
independent of that genus calculation and requires none of its geometry
as an additional axiom of the arithmetic proof.

\subsection{A sharper quartic calculation and a divisor algorithm}

The symmetric quartic family admits an especially transparent certificate
and sharper numerical bounds.  Let $a,b\in\mathbb Z$ and put
$\Delta=a^2-4b$.  Assume $\Delta\neq0$.  A solution of
\[
                         y^2=x^4+ax^2+b
\]
satisfies
\[
 (2x^2+a-2y)(2x^2+a+2y)=\Delta.
\tag{QF}
\]
Write the two factors as $u,v$ and let $d=|\Delta|$.
Both are nonzero and $|u|+|v|\leq d+1$ by
Lemma~\ref{lem:runge-product-bound}.  Since
$u+v=4x^2+2a$ and $v-u=4y$, we obtain the useful signed bounds
\[
                    4x^2\leq d+1-2a,
                 \qquad 4|y|\leq d+1.
\tag{QB}
\]
If the first right side is negative, the equation has no integral points.
Otherwise one may take
\[
 H=\left\lfloor\sqrt{\left\lfloor(d+1-2a)/4\right\rfloor}\right\rfloor,
 \qquad K=\left\lfloor(d+1)/4\right\rfloor.
\]
For the more general certificate theorem one could simply take
$D=2$, $A=2X^2+a$, and $R=4b-a^2$, giving the less sharp but still explicit
bound $|x|\leq |a|+|\Delta|$.

There is also a complete algorithm that avoids scanning this rectangle.
Enumerate the nonzero integers $u$ in $[-d,d]$ and retain those dividing
$\Delta$.  Set $v=\Delta/u$.  For each ordered factor pair test whether
\[
           z=\frac{u+v-2a}{4},\qquad y=\frac{v-u}{4}
\]
are integers, whether $z\geq0$, and whether $z$ is an integer square.
If $z=x_0^2$, return $(x_0,y)$ and $(-x_0,y)$, removing duplicates.
Every solution supplies exactly such a factor pair by (QF), so none is
missed.  Conversely the displayed reconstruction gives
$2z+a=(u+v)/2$ and $2y=(v-u)/2$; their difference of squares is
$uv=\Delta$.  Expanding and substituting $\Delta=a^2-4b$ yields
$y^2=z^2+az+b$, the original equation with $z=x_0^2$.  Thus every returned
pair is valid.  Divisibility, the congruences modulo $4$, and the square
test are all bounded decisions.

For a concrete positive-genus example, choose $a=7$ and $b=1$, so
$\Delta=45$.  The bounds give $x^2\leq8$ and $|y|\leq11$.
In (QF), $u+v=4x^2+14>0$ and $uv=45>0$, so both factors are positive.
Up to order the positive factor pairs are
$(1,45)$, $(3,15)$, and $(5,9)$.  They give respectively
$z=8,1,0$.  The first is not a square; the latter two give exactly
\[
                         (0,\pm1),\qquad(\pm1,\pm3).
\]
This is a complete list of six points.  It is verified by the factor
argument, not by interpreting a finite experimental search as complete.

More generally the quartic polynomial $X^4+aX^2+b$ is squarefree when
$b\Delta\neq0$.  A repeated root would have either $x=0$, forcing $b=0$,
or $2x^2+a=0$, forcing $b-a^2/4=0$.  Both alternatives are excluded.
Its smooth completion is therefore a genus-one curve, with two rational
points at infinity.  The assumptions hold for $a=7,b=1$.
The arithmetic bound itself only requires $\Delta\neq0$, so it also
handles some singular examples.  When $\Delta=0$, integrality forces $a$
to be even and
$X^4+aX^2+b=(X^2+a/2)^2$; in that case there are infinitely many integral
points.  The excluded zero remainder is therefore a necessary distinction
for this particular family.


\section{Why the elementary proofs fit iPRA}
\label{sec:arithmetic-audit}

This section explains the claimed weak-arithmetic status of the written
proofs. It does not identify native Lean induction with induction inside
iPRA. The intended object theory is the conventional first-order theory
implemented by ProveIt's \texttt{PRAAx}.

\subsection{Language, axioms, and rules}

The underlying domain is $\N$. There is one function symbol for each
primitive recursive construction: zero, successor, projections,
composition, and primitive recursion. For $f$ of arity $k$ and $g$ of
arity $k+2$, the recursion symbol obeys
\[
 \begin{split}
 h(0,\vec u)&=f(\vec u),\\
 h(n+1,\vec u)&=g(n,h(n,\vec u),\vec u).
 \end{split}
\]
Composition and projections have their usual defining equations.
The arithmetic axioms include $Sx\ne0$ and $Sx=Sy\to x=y$.
For each quantifier-free formula $\theta(n,\vec u)$, the induction
schema is
\[
 \theta(0,\vec u)\to
 \bigl[\forall n\,(\theta(n,\vec u)\to\theta(n+1,\vec u))\bigr]
 \to\forall n\,\theta(n,\vec u).
\]
The logic is intuitionistic first-order natural deduction with equality:
the introduction and elimination rules for $\wedge,\vee,\to,\forall,\exists$,
falsity elimination, equality reflexivity, and substitution of equals.
There is no unrestricted excluded middle or double-negation elimination.
Decidable arithmetic cases, such as the sign of an integer, are derived
using primitive recursive characteristic functions.

Signed integers use the canonical natural-number encoding in
Section~\ref{sec:logical-audit}. The integer ring laws, order laws,
absolute-value identities, and powers translate into equations and
inequalities between these codes. Their familiar elementary proofs use
only quantifier-free induction. Writing $x\in\Z$ in the arithmetic proofs
abbreviates this representation, not a new unbounded sort or quotient
principle added to the object theory.

\subsection{Concrete finite data}

Use dense coefficient arrays in increasing degree order. Normalize an
array by deleting trailing zero coefficients, with the empty array as
the zero polynomial. Length, entry lookup with default zero, truncation,
normalization, coefficient sums, addition, scalar multiplication, and
convolution are primitive recursive operations on finite sequence codes.
Horner evaluation is the recursion
\[
 E([],x)=0,\qquad E(a::L,x)=a+xE(L,x).
\]
Every polynomial identity in a certificate is checked over a specified
finite coefficient range. With array lengths $l_f,l_A,l_R$, it suffices
to test indices below $\max(l_f,2l_A,l_R)$. Outside that range both
sides are zero by their support bounds. This is why the certificate
predicate is primitive recursive rather than a universally quantified
semantic identity.

Rational arithmetic can be represented by a signed numerator and a
strictly positive denominator. Reduction to lowest terms is optional.
Addition and multiplication use the usual cross-multiplication formulas;
equality is checked by cross multiplication. If canonical rationals are
desired, the greatest common divisor can be computed by scanning divisors
up to the smaller nonzero input, with explicit finite fuel. None of the
proofs needs an unbounded search for a rational square root.

\subsection{Induction for coefficient estimates}

Fix a degree bound $m\ge1$, an integer $x$, and $t=|x|\ge1$.
For $k\le m$, define
\[
 T_k=\sum_{i<k}a_ix^i,\qquad S_k=\sum_{i<k}|a_i|.
\]
The required lower-coefficient estimate follows from the invariant
\[
 k\le m\quad\longrightarrow\quad |T_k|\le S_k t^{m-1}.
\]
The base case is $T_0=S_0=0$. At stage $k<m$, the triangle inequality
and $t^k\le t^{m-1}$ give
\[
 |T_{k+1}|\le |T_k|+|a_k|t^k
             \le(S_k+|a_k|)t^{m-1}.
\]
For $k>m$ the implication is inactive; alternatively, one freezes
the state after $m$ steps. All sums, powers, comparisons, and
absolute values are terms of primitive recursive functions. Thus the
displayed invariant is quantifier-free after definitions are expanded.
Monotonicity of powers for base $t\ge1$ follows by the same kind of
equational induction; no induction over an arbitrary family of formulas
is used.

The corresponding estimate for $R$ uses its computed leading index and
the same invariant. Its root bound is a finite sign/order argument at a
given $x$, not a search for all roots. The factor lemma uses only
\[
 |uv|=|u||v|,\qquad u\ne0\Rightarrow |u|\ge1,
 \qquad |s+t|\le|s|+|t|.
\]
When arguing by contradiction that $|x|\le B_x$, we use decidability
of that integer inequality: either it holds, or the explicit alternative
$|x|>B_x$ gives a contradiction. This is intuitionistically valid.

\subsection{Induction for square completion}

Let $u>0$ be the supplied integer square root of the leading coefficient
of $f$. The constructor has exactly $m$ descending coefficient stages
after the leading coefficient of $Q$ is set to $u$. Its state contains
a finite rational array $Q_k$, frozen after stage $m$. For $0\le k\le m$,
after $k$ stages the leading $k+1$ coefficients of $Q_k^2$ agree with
those of $f$. Encode this as a guarded bounded coefficient test
\[
 k\le m\quad\longrightarrow\quad
 \prod_{j=0}^{k}
  \operatorname{eq}_{\Q}\bigl([X^{2m-j}]Q_k^2,
                              [X^{2m-j}]f\bigr)=1.
\]
The product is an iterated Boolean conjunction. Its defining recursion
makes the assertion an atomic equality of primitive recursive terms.
The new coefficient occurs linearly with nonzero multiplier $2u$;
cross multiplication verifies the update. A second invariant states
that coefficients fixed in earlier stages are unchanged. Each step
therefore has a quantifier-free proof.

Representing denominators by products, rather than least common
multiples, avoids any issue about finding a minimum. The resulting
denominator $D$ and integer polynomial $A=DQ$ are primitive recursive
in the input array and the supplied positive square root $u$ of its
leading coefficient. Multiplying through the finite rational identities
proves $D^2f=A^2+R$. The test that the residual array is nonzero is
finite. If it is zero, the routine reports the square case and makes
no finiteness claim. Testing whether the leading coefficient is a
positive square itself requires only scanning up to that coefficient.

Uniform primitive recursiveness here refers to this \emph{fixed}
algorithm on variable input data. It is not a claim that one can
evaluate every coded primitive recursive program by a single primitive
recursive universal evaluator. There is no such general evaluator.
The implementation task is to construct one finite \texttt{PR} expression
for each fixed operation, then compose them.

\subsection{Exhaustive enumeration and its invariant}

Once $B_x$ is computed, one may scan $x=-B_x,\ldots,B_x$. For each
$n=f(x)\ge0$, compute the greatest $s\le n$ with $s^2\le n$ by a
bounded scan, or use a separately justified fast integer-square-root
algorithm. Retain $(x,s)$ and, for $s>0$, $(x,-s)$ precisely when
$s^2=n$. Every step has an explicit bound. This implementation avoids
the larger rectangular scan using $B_y$, while the latter remains a
simple proof-theoretic specification.

After $k$ horizontal positions have been processed, the retained list
contains exactly the solutions at those positions. In the rectangular
version, encode this as the conjunction over the first $k$ box indices
that each successful point test has its decoded tuple in the accumulated
list, and that every accumulated tuple passes the test and has a
processed index. All these quantifiers are bounded by the current code
lengths or the explicit box size. Their characteristic functions turn
the invariant into a quantifier-free formula. The single update either
appends the next successful tuple or leaves the list unchanged, which
proves the next invariant by finite membership identities.

At termination, the invariant proves exactness within the box. The
Runge theorem independently proves that every integral point lies in
the box. Combining these two conclusions proves global completeness.
The program is therefore an exhaustive solver on the stated certified
family, not just an experimental search below a chosen cutoff.

\subsection{What remains to be formalized}

The written arguments give the elementary arithmetic proof and its
allowed induction invariants. Completing the embedded derivation still
requires coding those functions as actual \texttt{PR} terms, deriving
their equations in \texttt{PRAAx}, constructing the characteristic
formulas and their \texttt{QF} witnesses, and assembling natural-deduction
proof terms. These are concrete finite tasks for the elementary family.
Their completion is not asserted by the presence of a native Lean
specification or by the Python checks.

The external identification of the examples' projective genus uses
standard algebraic geometry. The integral-point theorem itself does
not invoke that identification as an axiom: its hypotheses and
conclusion are directly arithmetic. A separate formalization of the
geometric interpretation would be needed to make the genus statement
part of the same internal arithmetic development.

\section{The general theorem: a precise remaining route}
\label{sec:general-route}

The elementary proof leaves the original general request open. This section
identifies a constructive target that does not demand a height bound, proves
its finite combinatorial last step, and specifies the missing dependencies.
It is a research program, not a claim that its hard theorem has been proved.

\subsection{A short audit of the classical route}

The Subspace Theorem approach uses a Riemann--Roch space of functions with
controlled poles at the boundary. Local estimates turn integral points
into solutions of linear-form inequalities. The Subspace Theorem confines
their images to finitely many proper subspaces; a nonzero function with
bounded poles has only boundedly many zeros, so each subspace contributes
finitely many curve points. Covers and lifting arguments handle the
positive-genus cases with too few boundary points
\cite{CorvajaZannier2002}.

This outline has not verified the strength of those dependencies in iPRA.
Qualitative compactness, arbitrary infinite subsequences, an unlocated
exceptional subspace, or a nonconstructive cover cannot be accepted as
hidden axioms. Each use needs a finite replacement with explicitly coded
data and permitted induction.

Quantitative Subspace Theorems control the \emph{number} of exceptional
subspaces under explicit hypotheses; they do not thereby construct them.
Evertse's Theorems 2.1--2.2 give relevant quantitative forms
\cite{Evertse2010}. The following finite formulation avoids assuming that
a global exceptional set has already been found.

\subsection{A finite covering lemma}

\begin{proposition}[Finite cover to cardinality]
Let $P$ be a decidable predicate on canonical point codes, let $\Phi$ map
point codes to finite vectors over a coded field, and let $M,D\in\N$.
Suppose the following two contracts have been proved.

For every finite list $p_1,\ldots,p_t$ of pairwise distinct points satisfying
$P$, the vectors $\Phi(p_i)$ admit a cover by at most $M$ proper linear
subspaces, with decidable membership in each supplied subspace.

For every subspace $W$ occurring in such a cover, at most $D$ distinct
points in the list have image in $W$.

Then no list of $MD+1$ pairwise distinct points satisfying $P$ exists.
\end{proposition}

\begin{proof}
Suppose $t=MD+1$ distinct points are given. Obtain a cover
$W_1,\ldots,W_s$ with $s\le M$. For each of the finitely many points,
test the subspaces in order and assign it to the least index $j$ for
which $\Phi(p_i)\in W_j$. The covering contract ensures that the finite
search succeeds. No selection from an infinite family occurs.

Let $n_j$ count the points assigned to $j$. Each point is assigned
exactly once, so $t=\sum_{j=1}^s n_j$. Every point counted by $n_j$
lies above $W_j$, and the input points are distinct; hence $n_j\le D$.
Consequently $MD+1=t\le sD\le MD$, a contradiction. These sums,
least indices, and counts are primitive recursive, and their
equalities and inequalities have quantifier-free finite-loop proofs.
\end{proof}

Using the proposition for curves is conditional on proving its two
contracts. It returns no height bound. If there are several local
inequality systems, with respective bounds $M_i,D_i$, and at most $E$
points in an explicitly enumerated exception box, the same assignment
argument gives the cardinality bound
\[
                         E+\sum_i M_iD_i.
\]
No point outside the exception box has thereby been located.

\subsection{Why finite covers can have finite certificates}

Fix a list of $t$ vectors in dimension $n$. If a proper subspace $W$
covers a sublist, replace $W$ by the span of that sublist. This new
space remains inside $W$ and is proper. Finite Gaussian elimination
selects a basis of at most $n-1$ input vectors, including the empty
basis for a zero-dimensional span.

For a particular tuple, therefore, a covering certificate may be sought
among lists of at most $M$ subsets of at most $n-1$ of its $t$ vectors.
There are at most $2^t$ candidate subsets and at most $(2^t+1)^M$
padded lists of candidates. Rank, dependence, and membership have
finite exact checks over rational numbers, or over a supplied
number-field basis. These search limits are primitive recursive.

The existence of a successful certificate for \emph{every} relevant tuple
is still a mathematical theorem. The search can fail; executing it does
not show that it always succeeds. The needed statement is a finite-tuple
version of the quantitative Subspace Theorem, proved with the induction
permitted in iPRA. No such proof is provided here. Assuming a global
finite set of exceptional subspaces would silently introduce stronger
positive information.

\subsection{The geometric input to the fibre bound}

Suppose $\phi_1,\ldots,\phi_n$ are linearly independent rational
functions on a smooth projective curve whose poles are all bounded
by an effective divisor $E$ of degree $D$. A proper linear
subspace admits a nonzero annihilating linear functional. Its
coefficients give a nonzero function $\sum a_i\phi_i$, by linear
independence. Its poles are bounded by $E$. Equality of the degrees
of its zero and pole divisors bounds its distinct zeros by $D$.

This conditional geometric argument suggests the second contract.
It does not construct the functions or their divisor data in arithmetic.
An internal development must represent the function field, divisor,
vector space, and independence certificate, then prove the zero-count
lemma from those finite descriptions. An abstract divisor theorem in
the host Lean library does not automatically discharge that obligation.

\subsection{Explicit unresolved obligations}

A general cardinality proof still needs finite presentations of the
Riemann--Roch functions and their local expansions; a proof of the
finite-tuple quantitative approximation theorem; bounds for its
auxiliary polynomials and matrices using permitted induction; control
of local case splits and the small-height part; and a quantitative
cover-and-lift construction for the positive-genus one- or two-boundary
cases, including all field and denominator changes. Finally, one must
prove the zero-count contract and apply the finite covering lemma.

These are mathematical dependencies. Describing them as standard,
declaring Lean axioms for them, or storing them as unproved structure
fields would not complete the requested proof.

For the positive height/list version, the target is stronger: a
uniform function bounding coordinates in the original affine model,
with a correctness proof. The Runge theorem supplies this for its
explicit family. It does not supply the general bound discussed by
the primary sources. No independence or unprovability conclusion is
drawn from the unresolved state of that step.


\section{Lean definitions, checked results, and completion plan}
\label{sec:lean}

\subsection{Two different formalization levels}

The sources use Lean 4.32.0 and its bundled \texttt{Std} library.
They contain native definitions and kernel-checked arithmetic theorems,
not yet proof objects in ProveIt's first-order theory. Lean's metatheory
is stronger than iPRA even when the mathematical argument has a short
iPRA proof.

The arithmetic and search modules are \texttt{Arithmetic.lean} and
\texttt{FiniteSearch.lean}. The other modules are
\texttt{QuarticExample.lean} and \texttt{Specifications.lean}.
All four compile.
The audit module prints axiom dependencies and checks target types.
The general targets are definitions of propositions, not asserted
theorems or added axioms.

\subsection{Checked native mathematics}

\texttt{Arithmetic.lean} proves the nonzero-product bound, the
difference-of-squares identity, and the decisive gap statement
\[
 z^2=a^2+r,\quad r\ne0\quad\Longrightarrow\quad |a|\le|r|.
\]
It proves the quartic factorization and both sharp quartic bounds,
uniformly in $a,b,x,y$, and proves that no natural-number square lies
strictly between consecutive natural-number squares.

\texttt{FiniteSearch.lean} defines
\begin{lstlisting}
def scan (p : Nat -> Bool) (B : Nat) : List Nat :=
  (List.range (B+1)).filter p
\end{lstlisting}
and proves its exact membership characterization. It proves that a
complete finite list and a numerical bound are equivalent for predicates
on natural-number codes. It also proves the intentionally weaker
fact that every finite prefix is bounded.

\texttt{QuarticExample.lean} proves for all integers $x,y$ that
\[
\begin{split}
y^2=x^4+7x^2+1\quad\Longleftrightarrow\quad
&[x=0\ \wedge\ (y=1\vee y=-1)]\\
{}\vee{}&[(x=1\vee x=-1)\ \wedge\ (y=3\vee y=-3)].
\end{split}
\]
Its global coordinate bounds are proved before the finite case analysis.
This is not a bounded experiment promoted to a universal theorem.

The axiom audit contains only the standard Lean principles
\texttt{propext} and \texttt{Quot.sound}; the finite-prefix witness
requires none. No audited result depends on \texttt{Classical.choice},
\texttt{sorryAx}, or an added theorem axiom. This fact does not replace
the missing embedded iPRA derivations.

\subsection{Exact specifications and unproved targets}

The specification file defines polynomials as \texttt{List Int},
Horner evaluation, coefficient lookup with default zero, a coefficient
norm, and convolution coefficients. A certificate contains the finite
data $f,A,R,D$. The Boolean \texttt{certificateCheck} verifies positivity,
degree, nonzero conditions, and coefficient equality over a finite
support range. It contains no assumed height bound.

The bound functions are implemented directly:
\[
\begin{split}
\texttt{xBound}(c)&=\|A_{<\deg A}\|_1+\|R\|_1,\\
\texttt{yBound}(c)&=\sum_i |f_i|\max(1,\texttt{xBound}(c))^i.
\end{split}
\]
The file also defines the integer interval, rectangular filter, and
$X^{2m}+X+1$. The propositions \texttt{RungeBoundTarget},
\texttt{RungeEnumerationTarget}, and \texttt{GenusFamilyTarget}
are unproved in Lean. The corresponding mathematics is proved in
the paper. A concrete sample certificate is checked, and its finite
box evaluates to the expected four points.

For general affine systems the file defines monomials, exponent lists,
polynomial evaluation, a well-formedness test, an integral-point test,
and coordinate height. \texttt{PositiveSiegelTarget} and
\texttt{CountSiegelTarget} have different conclusions. Each accepts an
\texttt{admissible} predicate as a parameter. Implementing that
predicate so that it really certifies smooth positive-genus curves,
and proving its coverage, remain separate obligations. A parameterized
definition asserts no geometric theorem.

\subsection{The actual ProveIt interfaces}

At the pinned revision the namespace is
\texttt{LeanProofs.ConciseIPRA}. The inspected interfaces are
\begin{center}
\begin{tabular}{@{}p{0.35\linewidth}p{0.56\linewidth}@{}}
\toprule
Interface & Role\\
\midrule
\texttt{PR k} & Primitive recursive construction of arity $k$\\
\texttt{PTm n} & Arithmetic terms with $n$ free variables\\
\texttt{Fm PTm n} & First-order formulas\\
\texttt{Fm.QF} & Quantifier-free restriction\\
\texttt{PRAAx} & Arithmetic axioms and restricted induction\\
\texttt{Prf PRAAx Gamma phi} & Intuitionistic natural-deduction derivability\\
\texttt{PRA.qfInd} & Quantifier-free induction rule\\
\texttt{pra\_iff\_native} & Transfer to the concise native calculus\\
\bottomrule
\end{tabular}
\end{center}

The intended elementary result must ultimately have the type
\begin{lstlisting}
Prf PRAAx [] rungeFormula
\end{lstlisting}
where \texttt{rungeFormula} is the closed universal formula built
from the compiled \texttt{PR} functions and actual \texttt{Fm}
constructors. A host proposition with similar mathematical words
is insufficient. Once the derivation exists,
\texttt{pra\_iff\_native} transports it; the equivalence does not
create the missing derivation.

\subsection{Dependency-ordered completion tasks}

First prove that the finite polynomial checker implies the evaluated
identity. Establish the coefficient and integer-root bounds for
normalized lists. Combine them with \texttt{square\_gap} to prove
\texttt{RungeBoundTarget}. Prove membership in the integer interval
and nested filter to obtain \texttt{RungeEnumerationTarget}.
Instantiate the certificate uniformly in $m$ and apply the consecutive
square lemma to prove \texttt{GenusFamilyTarget}. Every mathematical
ingredient appears in Section~\ref{sec:runge-certificate}.

Next construct fixed primitive recursive codes for signed arithmetic,
arrays, sums, evaluation, the certificate test, both bounds, and
enumeration. Derive their equations inside \texttt{PRAAx}. Compile
the invariants of Section~\ref{sec:arithmetic-audit} to characteristic
equations, supply their \texttt{QF} witnesses, and use
\texttt{PRA.qfInd}. Assemble the universal derivation by natural
deduction and equality substitution, then specialize it to concrete
certificates and transport it to the concise calculus.

The accompanying \texttt{OBJECT\_PLAN.txt} expands these steps into
individual invariants and proof obligations. Completion requires
the actual object derivation, an axiom audit, and a clean build.
A native Lean build alone does not meet that requirement.

There is a finite implementation path for the elementary results.
The general Siegel dependencies in Section~\ref{sec:general-route}
are additional mathematical research tasks. Completing routine Lean
infrastructure is not promised to resolve them.


\section{Verification, reproducibility, and remaining scope}
\label{sec:verification}

\subsection{What was checked}

The exact reference code was run with Python 3 using only standard-library
integer and rational arithmetic. The results recorded in
\texttt{code/verification.json} include:
\begin{center}
\begin{tabular}{@{}p{0.68\linewidth}r@{}}
\toprule
Check & Cases\\
\midrule
Bounded binary square root versus an independent implementation & 10,385\\
Family $X^{2m}+X+1$, $m=2,\ldots,12$ & 11\\
Quartics $X^4+aX^2+b$, $a,b\in[-8,8]$, $\Delta\ne0$ & 284\\
Additional square-completion inputs & 2\\
Malformed or out-of-scope inputs rejected & 9\\
\bottomrule
\end{tabular}
\end{center}

For quartics, the certificate enumerator, signed-divisor enumerator,
and an independent direct-power search agree. The independent search
uses a larger horizontal window and Python's separate integer-square-root
routine. These are implementation checks. They are not the proof
that the window contains every point; that fact comes from the
universal arithmetic bound. The family theorem for every $m\ge2$ is
likewise proved symbolically, not inferred from its eleven tested values.

The binary square-root implementation has explicit fuel. For $n\ge0$
it starts at $\ell=0,h=n+1$, maintains
$\ell^2\le n<h^2$, and replaces an endpoint by the integer midpoint
according to a square comparison. The interval length is at most
halved up to rounding. For $n\ge1$ with bitlength $b$, initially
$h-\ell=n+1\le2^b$, so $b+1$ iterations suffice to reach adjacent
endpoints. The case $n=0$ begins with adjacent endpoints. This proves
termination with a primitive recursive bound; it is not an
unbounded perfect-square search.

\begin{samepage}
The Lean build used version 4.32.0, commit
\begin{center}\small\ttfamily
8c9756b28d64dab099da31a4c09229a9e6a2ef35.
\end{center}
\end{samepage}
Its reproducible output is in \texttt{verification/lean\_build.txt}.
The exact theorem names and axiom dependencies are printed there.
A compiled definition of an unproved target is reported as a
specification, not as a proof. No complete ProveIt repository build
was performed.

The mathematical arguments were separately reviewed for their
coefficient estimates, nonzero remainder requirement, square
completion, sign cases, geometric interpretation, and logical
witness claims. The coding discussion was clarified to require
canonical point codes. The symbol $\Delta=a^2-4b$ is the auxiliary
quartic parameter; it is not the full polynomial discriminant.
The full quartic discriminant is $16b\Delta^2$.

\subsection{Archive contents}

\begin{center}
\begin{tabular}{@{}p{0.37\linewidth}p{0.54\linewidth}@{}}
\toprule
Path & Purpose\\
\midrule
\texttt{paper/siegel\_ipra.pdf} & Complete report\\
\texttt{paper/*.tex} & Master source and included sections\\
\texttt{lean/*.lean} & Checked native modules, specifications, and audit\\
\texttt{lean/OBJECT\_PLAN.txt} & Detailed object-iPRA compilation obligations\\
\texttt{lean/build.sh} & Reproducible Lean build\\
\texttt{code/runge.py} & Exact certificate constructor and solvers\\
\texttt{code/verify.py} & Independent finite regression checks\\
\texttt{code/build\_pdf.py} & Ordinary or private-format LaTeX build\\
\texttt{verification/} & Build and document-validation records\\
\texttt{provenance/} & Pinned source audit and research references\\
\texttt{README.txt} & Entry point, scope, and reproduction commands\\
\texttt{SHA256SUMS} & Checksums of the delivered files\\
\bottomrule
\end{tabular}
\end{center}

From the extracted root directory, run:
\begin{lstlisting}
python3 code/verify.py
bash lean/build.sh
python3 code/build_pdf.py
\end{lstlisting}
Use the version pinned in \texttt{lean/lean-toolchain}. If Lean is
available at a custom path, set \texttt{SIEGEL\_LEAN\_BIN} to that
executable before running the build script. The PDF script uses
ordinary \texttt{pdflatex} when configured and can prepare a private
format from installed TeX sources when the system format is missing.
It runs without shell escape.

Verify \texttt{SHA256SUMS} before rebuilding, since regenerated
logs and metadata may change. The archive contains source and
evidence; it does not bundle a Lean distribution or claim to have
recompiled the upstream foundation.

\subsection{The exact unresolved conclusion}

The fully developed arithmetic theorem proves effective finiteness
for the certified Runge family and supplies actual complete lists.
The archive also contains compiled native Lean proofs of important
lemmas and one global positive-genus example. The object-iPRA
derivations remain to be implemented using the explicit induction
invariants.

For arbitrary smooth positive-genus affine curves over number fields,
the requested general iPRA proof is not provided. A positive complete
list or height-bound proof would require the missing general effective
argument. A cardinality proof can pursue the different finite-tuple
route discussed above, whose approximation and geometry dependencies
have not been established in iPRA here. This report does not prove
either general target, and it does not claim either one impossible.


\begingroup\small
\begin{thebibliography}{99}

\bibitem{Salehi2002}
S. Salehi,
\emph{Decidable Formulas of Intuitionistic Primitive Recursive Arithmetic},
Reports on Mathematical Logic \textbf{36} (2002), 55--61.
Fact 3.5, p.~57, states the arbitrary-matrix primitive recursive
witness theorem used here.
\url{https://rml.tcs.uj.edu.pl/rml-36/36-SALEHI.pdf}.

\bibitem{Wehmeier1997}
K. F. Wehmeier,
\emph{Fragments of HA based on $\Sigma_1$-induction},
Archive for Mathematical Logic \textbf{37} (1997), 37--49.
\url{https://doi.org/10.1007/s001530050081}.
The publisher abstract and metadata were checked; the witness formulation
used in this report is the explicit statement in Salehi.

\bibitem{CorvajaLombardoZannier2024}
P. Corvaja, D. Lombardo, and U. Zannier,
\emph{Examples of effectivity for integral points on certain curves of genus 2},
arXiv:2411.17930v2, revised 28 December 2024.
Introduction and Remark 1.2.
\url{https://arxiv.org/abs/2411.17930};
\url{https://arxiv.org/pdf/2411.17930}.

\bibitem{BiluIllengo2009}
Y. Bilu and M. Illengo,
\emph{Effective Siegel's theorem for modular curves},
arXiv:0905.0418 (2009), Introduction and Theorem 1.2.
\url{https://arxiv.org/pdf/0905.0418}.

\bibitem{CorvajaZannier2002}
P. Corvaja and U. Zannier,
\emph{A subspace theorem approach to integral points on curves},
Comptes Rendus de l'Acad\'emie des Sciences, Paris, S\'erie I
\textbf{334} (2002), 267--271.
\url{https://doi.org/10.1016/S1631-073X(02)02240-9}.
Open text:
\url{https://www.numdam.org/item/10.1016/S1631-073X(02)02240-9.pdf}.

\bibitem{Evertse2010}
J.-H. Evertse,
\emph{On the quantitative Subspace Theorem},
arXiv:1008.2268 (2010), especially Theorems 2.1--2.2.
\url{https://arxiv.org/pdf/1008.2268}.
Used to identify the quantitative dependency, not as an iPRA formalization.

\bibitem{Levin2016}
A. Levin,
\emph{Integral points and the Shafarevich conjecture for abelian surfaces},
in \emph{Diophantische Approximationen},
Oberwolfach Report 21/2016, 1176--1178, especially 1176--1177.
\url{https://doi.org/10.4171/OWR/2016/21}.
This source further explains the effectivity distinction for genus two
with different choices of poles.

\bibitem{ProveIt2026}
V. Reshetnikov and contributors,
\emph{ProveIt}, repository snapshot:
\begin{quote}\small\ttfamily
7eeb3c0e821594d5d5ebcfbaf4bdba64819e3502.
\end{quote}
\url{https://github.com/VladimirReshetnikov/ProveIt}.
Inspected foundation files:
\begin{quote}\small\ttfamily
Logic/PrimitiveRecursiveArithmetic/ConciseIPRA/Lean/\\
ConciseIPRA/PR.lean\\
ConciseIPRA/Deduction.lean\\
ConciseIPRA/NativeBiInterpretation.lean
\end{quote}
Related packet index and status files are recorded with their exact
paths and hashes in \texttt{provenance/repository\_manifest.json}.
No upstream modification or full upstream build is asserted.

\bibitem{UserReference}
\emph{Siegel's theorem on integral points}, Wikipedia,
the reference supplied with the request; accessed 15 November 2026.
\url{https://en.wikipedia.org/wiki/Siegel\%27s_theorem_on_integral_points}.
The primary sources above support the mathematical and logical
status discussion.

\end{thebibliography}
\endgroup


\end{document}
